\begin{textsample}{POS Dim 1 – rj – Score 23.00 – rj/rj\_ef\_ciencias\_humanas\_his\_2.txt}  \label{ex:f1_pos_258}
A História como Campo Científico : Perspectivas de \textbf{Abordagem} e Questões \textbf{Metodológicas} A história se \textbf{configurou} como \textbf{um} campo autônomo do conhecimento apenas no \textbf{século} XIX, quando, \textbf{influenciada} pelo pensamento positivista procurou se colocar como \textbf{uma} ciência nova, com metodologia própria e bem definida.
Daí em diante, passou a ser componente curricular definitivo nos sistemas de ensino por todo mundo.
Obviamente, as perspectivas \textbf{conceituais} e \textbf{metodológicas} desse campo do saber se ampliaram desde \textbf{século} XIX até esse início do \textbf{século} XXI.
Novas correntes surgiram na historiografia ao final do \textbf{século} XIX e ao longo do \textbf{século} XX, e \textbf{operaram} a crítica a visão positivista reinante no início.
Do Marxismo a História Cultural, da Escola dos Annales ( em suas diferentes gerações ) a História Social Inglesa, passando por outras correntes historiográficas, a existência de \textbf{uma} em multiplicidade de correntes fez o nosso campo se desenvolver, tentando responder aos desafios e as questões de cada época, e sempre se articulando a perspectivas \textbf{teóricas} de outras áreas.
A emergência de novos \textbf{atores} sociais ( afrodescendentes, mulheres, indígenas, a questão das “ minorias ” e da diversidade, etc... ) refletiu na historiografia, \textbf{que} não pode mais negar a existência desses setores, nem de suas pautas.
Isso se refletiu também no ensino da \textbf{disciplina}, \textbf{que} teve \textbf{que} incorporar esses setores na sala de aula ( até porque eles estão lá representados também ).
Nesse contexto, a história assume o papel de conhecimento indispensável a toda a área das ciências humanas, pois é ela \textbf{que} vai fornecer a noção de tempo, \textbf{central} para entendermos como os seres humanos encontraram respostas as questões \textbf{que} foram \textbf{centrais} em diferentes momentos de sua existência, e, da mesma forma, ao refletirmos sobre isso, estamos refletindo sobre nosso próprio presente e nossa própria existência.
Noções como tempo histórico, processo ( longa, média ou curta duração ) ruptura/transformação, e continuidade são ferramentas \textbf{centrais} a todo conhecimento histórico.
Com a história, podemos perceber \textbf{que} cada indivíduo é “ produto ” de seu tempo, mas também é fruto de \textbf{uma} gama de experiências, saberes e condições produzidas no passado, sendo \textbf{que} essas experiências, saberes e condições também definem o seu lugar no presente.
Fazer essa reflexão é essencial para a construção do saber histórico.
No Brasil, nossa \textbf{disciplina} passou a fazer parte dos currículos escolares a partir de 1827.
A exaltação aos “ heróis e mitos ” fundadores de \textbf{uma} nação idealizada ( ou “ inventada ”, como \textbf{diriam} muitos ), era a tônica das obras históricas do período, bem como a preocupação principal do ensino dessa \textbf{disciplina}.
A história aparecia, então, dentro de \textbf{uma} perspectiva dos “ vencedores ”, ignorando quase \textbf{que} por completo os setores \textbf{que} não faziam parte de \textbf{uma} determinada “ elite ” dominante, excluindo e marginalizando o \textbf{que} não cabia dentro de \textbf{um} paradigma eurocêntrico e cientificista da sociedade.
Não à toa, Varnhagen, foi a principal referência desse período.
No final do \textbf{século} XIX, início do \textbf{século} XX começa a desconstrução dessa visão, com a contribuição de historiadores como Caio Prado Junior, Sérgio Buarque de Holanda e Gilberto Freyre ( antes havia existido Capistrano de Abreu ), \textbf{influenciados} pelas novas correntes da historiografia europeia e pelo marxismo.
Entretanto, foi somente após o fim do regime militar em 1985, \textbf{que} \textbf{uma} \textbf{abordagem} \textbf{metodológica} \textbf{que} incluía os “ vencidos ”, os “ excluídos ”, e outras \textbf{categorias} não representadas pelas visões mais tradicionais, começou a ganhar força na historiografia brasileira, dialogando com correntes como a Nova História e a História Social.
Porém, apesar dos estudos mais recentes \textbf{que} apontam para \textbf{uma} maior preocupação com outros \textbf{sujeitos} históricos, apesar da emergência de novas temáticas e recortes temporais, e apesar da relevância \textbf{que} o local passou a ter em detrimento do “ macro ”, essas tendências ainda não tinham se \textbf{cristalizado} na construção de propostas curriculares \textbf{que} agregassem essas novas \textbf{abordagens}.
Assim, o ensino de história na área de humanas deve reforçar a compreensão de \textbf{que} tudo é passível ( e possível ) de ser mudado, e \textbf{que}, é a ação consciente de \textbf{homens} e mulheres, \textbf{que} torna a mudança \textbf{uma} realidade, mas também deve reforçar a compreensão de \textbf{que} não haveria o \textbf{hoje}, se o ontem não tivesse existido da forma como existiu.
Mais \textbf{uma} vez, enfatizamos \textbf{que} essa relação passado/presente só se faz possível se \textbf{conseguirmos} desenvolver \textbf{uma} “ atitude historiadora ” em cada estudante, se cada indivíduo for capaz de comparar, interpretar, analisar e contextualizar o seu passado, com autonomia de pensamento, e pensamento crítico, para \textbf{que} o conhecimento histórico possa se construir de forma \textbf{sólida} e conectada com a realidade, para \textbf{que} assim, o ensino de história possa estar em consonância com o “ \textbf{espírito} ” da BNCC e da área de humanas, contribuindo, antes de tudo para a formação de seres \textbf{humanos} com valores éticos e dispostos a construir \textbf{uma} nova sociedade.

% matched lemmas: abordagem, ator, categoria, central, conceitual, configurar, conseguir, cristalizar, disciplina, dizer, espírito, hoje, homem, humanos, influenciar, metodológico, operar, que, sujeito, século, sólido, teórico, um
\end{textsample}
